\section{What The Boundary Carries Forward}

The chapter has not added a new primitive to the grammar. It has protected
the forced continuum reading from two temptations: unearned completion and
interior mixing. The extensional ledger identified duplicate presentations
without losing finite support. Countability survived because the quotient
only removed redundant finite indices. The continuum therefore remained
answerable to finite conditions.

The chapter then made order readable. The gauge commutator residue closed
only when successor counting could carry the ordered transports and their
comparison. That residue entered completion through the normed
Sobolev--Hilbert door. The finite origin was traced inside the completed
reading, and the geometry/gauge split lifted rather than dissolved.

Compatibility then did its most severe work. It forbade the completed
interior from forming a scalar mixed term between the geometric and gauge
channels. The term vanished not because the channels had no relation, but
because their relation was generated at the boundary and could not be
reinserted as an interior product without undoing the split.

Finally, the boundary trace read the terminal residue. Radiation was forced
there because the interior had no licensed mixed certificate left. The
boundary restored a place where the completed residue could be read without
violating the finite grammar that produced it.

The next chapter receives this boundary as an outside reader. Once radiation
has been forced to the boundary, the interior invariant can be read from
outside as an obstruction rather than as an interior possession. The
continuum has been completed, but completion has not silenced the ledger.
The boundary still carries the residue the interior cannot spend.
